\documentclass[10pt]{article}
\usepackage[utf8]{inputenc}
\usepackage[spanish]{babel}
\usepackage{amsmath, amssymb}
\usepackage{geometry}
\geometry{a4paper, margin=1.5cm}
\usepackage{multicol}
\usepackage{enumitem}
\usepackage{fancyhdr}
\usepackage{eso-pic}
\usepackage{graphicx}

% Configuración de página y encabezados/pies
\pagestyle{fancy}
\fancyhf{}
\fancyfoot[R]{Christian Cristobal Carbajal García}
\renewcommand{\headrulewidth}{1pt}
\renewcommand{\footrulewidth}{1pt}

% Logo y marca de agua
\makeatletter
\AddToShipoutPictureFG*{%
  \AtPageUpperLeft{%
    \put(20,-60){%
      \IfFileExists{/home/christian/Documentos/Reportes_Aut/imagenes/logo_pdf.png}{%
        
\includegraphics[width=2.5cm]{/home/christian/Documentos/Reportes_Aut/imagenes/logo_pdf.png}%
      }{}%
    }%
  }%
  \AtPageUpperLeft{%
    \put(470,-85){%
      \IfFileExists{/home/christian/Documentos/Reportes_Aut/imagenes/marca_de_agua.png}{%
        \includegraphics[width=4.0cm]{/home/christian/Documentos/Reportes_Aut/imagenes/marca_de_agua.png}%
      }{}%
    }%
  }%
}
\makeatother

\title{\textbf{Guía de Álgebra y Geometría Analítica -- 50 preguntas}}
\author{}
\date{}

\begin{document}

\maketitle

\noindent
\textbf{Instrucciones:} Selecciona la opción correcta en cada pregunta. Las respuestas numéricas pueden variar ligeramente en los decimales.

\begin{multicols}{2}
\raggedcolumns

\begin{enumerate}

% ÁLGEBRA (25 preguntas)

\item Resuelve la ecuación $3x - 7 = 2x + 5$. ¿Cuál es el valor de $x$?
\begin{enumerate}[label=(\alph*)]
\item $10$
\item $12$
\item $8$
\item $14$
\end{enumerate}

\item Si $x^2 - 5x + 6 = 0$, ¿cuáles son sus raíces?
\begin{enumerate}[label=(\alph*)]
\item $2$ y $3$
\item $-2$ y $-3$
\item $1$ y $6$
\item $-1$ y $-6$
\end{enumerate}

\item Simplifica $(2a+3b)^2$:
\begin{enumerate}[label=(\alph*)]
\item $4a^2+12ab+9b^2$
\item $4a^2+6ab+9b^2$
\item $2a^2+12ab+3b^2$
\item $4a^2+9b^2$
\end{enumerate}

\item Factoriza $x^2 - 16$:
\begin{enumerate}[label=(\alph*)]
\item $(x-4)(x+4)$
\item $(x-8)(x+2)$
\item $(x-4)^2$
\item $(x-16)(x+1)$
\end{enumerate}

\item Resuelve el sistema: $2x + y = 7$, $x - y = 2$:
\begin{enumerate}[label=(\alph*)]
\item $x=3$, $y=1$
\item $x=2$, $y=3$
\item $x=4$, $y=-1$
\item $x=1$, $y=5$
\end{enumerate}

\item ¿Cuál es el valor de $\log_2 8$?
\begin{enumerate}[label=(\alph*)]
\item $2$
\item $3$
\item $4$
\item $1$
\end{enumerate}

\item Despeja $x$ en $a x + b = c$:
\begin{enumerate}[label=(\alph*)]
\item $x = \frac{c-b}{a}$
\item $x = \frac{c+b}{a}$
\item $x = c - b - a$
\item $x = \frac{a}{c-b}$
\end{enumerate}

\item El producto $(x-2)(x^2+2x+4)$ es igual a:
\begin{enumerate}[label=(\alph*)]
\item $x^3 - 8$
\item $x^3 + 8$
\item $x^3 - 4x^2 + 8$
\item $x^3 + 4x^2 - 8$
\end{enumerate}

\item ¿Cuál es el valor de $2^{-3}$?
\begin{enumerate}[label=(\alph*)]
\item $\frac{1}{8}$
\item $8$
\item $\frac{1}{6}$
\item $6$
\end{enumerate}

\item La solución de $|x-3| = 5$ es:
\begin{enumerate}[label=(\alph*)]
\item $x = 8$ o $x = -2$
\item $x = 8$ o $x = 2$
\item $x = -8$ o $x = 2$
\item $x = -8$ o $x = -2$
\end{enumerate}

\item Si $\frac{x}{2} + 3 = 7$, entonces $x$ vale:
\begin{enumerate}[label=(\alph*)]
\item $8$
\item $5$
\item $10$
\item $4$
\end{enumerate}

\item Factoriza $3x^2 - 12$:
\begin{enumerate}[label=(\alph*)]
\item $3(x-2)(x+2)$
\item $3(x-4)(x+1)$
\item $3(x-2)^2$
\item $(3x-4)(x+3)$
\end{enumerate}

\item Resuelve $2^{x+1} = 16$:
\begin{enumerate}[label=(\alph*)]
\item $3$
\item $4$
\item $2$
\item $5$
\end{enumerate}

\item La suma de los primeros 10 números naturales es:
\begin{enumerate}[label=(\alph*)]
\item $55$
\item $50$
\item $45$
\item $60$
\end{enumerate}

\item Si $f(x) = 3x - 2$, entonces $f(5)$ es:
\begin{enumerate}[label=(\alph*)]
\item $13$
\item $15$
\item $10$
\item $17$
\end{enumerate}

\item El valor de $(2^3)^2$ es:
\begin{enumerate}[label=(\alph*)]
\item $64$
\item $32$
\item $16$
\item $8$
\end{enumerate}

\item Resuelve $x^2 - 4x + 4 = 0$:
\begin{enumerate}[label=(\alph*)]
\item $x=2$ (doble)
\item $x= -2$ (doble)
\item $x=4$ (doble)
\item $x=0$ y $x=4$
\end{enumerate}

\item Simplifica $\frac{x^2-9}{x-3}$ (para $x \neq 3$):
\begin{enumerate}[label=(\alph*)]
\item $x+3$
\item $x-3$
\item $x+9$
\item $x-9$
\end{enumerate}

\item Si $\log x = 2$, entonces $x$ es:
\begin{enumerate}[label=(\alph*)]
\item $100$
\item $20$
\item $10$
\item $2$
\end{enumerate}

\item El mínimo común múltiplo de $6$ y $8$ es:
\begin{enumerate}[label=(\alph*)]
\item $24$
\item $48$
\item $12$
\item $16$
\end{enumerate}

\item Resuelve $3x^2 + 5x - 2 = 0$:
\begin{enumerate}[label=(\alph*)]
\item $x = \frac{1}{3}$ y $x = -2$
\item $x = -\frac{1}{3}$ y $x = 2$
\item $x = 2$ y $x = -\frac{1}{2}$
\item $x = -2$ y $x = \frac{1}{2}$
\end{enumerate}

\item Si $a = 2$, $b = 3$, entonces $a^2 + b^2$ es:
\begin{enumerate}[label=(\alph*)]
\item $13$
\item $10$
\item $5$
\item $25$
\end{enumerate}

\item La ecuación de la recta con pendiente $2$ que pasa por $(1,3)$ es:
\begin{enumerate}[label=(\alph*)]
\item $y = 2x + 1$
\item $y = 2x + 3$
\item $y = 2x - 1$
\item $y = 2x - 3$
\end{enumerate}

\item El valor de $\sqrt{50}$ simplificado es:
\begin{enumerate}[label=(\alph*)]
\item $5\sqrt{2}$
\item $2\sqrt{5}$
\item $10\sqrt{5}$
\item $25\sqrt{2}$
\end{enumerate}

\item Si $3x - 2y = 5$ y $x + y = 2$, entonces $y$ vale:
\begin{enumerate}[label=(\alph*)]
\item $\frac{1}{5}$
\item $\frac{2}{5}$
\item $\frac{3}{5}$
\item $1$
\end{enumerate}

% GEOMETRÍA ANALÍTICA (25 preguntas)

\item La distancia entre los puntos $A(1,2)$ y $B(4,6)$ es:
\begin{enumerate}[label=(\alph*)]
\item $5$
\item $4$
\item $6$
\item $7$
\end{enumerate}

\item El punto medio del segmento con extremos $(-2,3)$ y $(4,-1)$ es:
\begin{enumerate}[label=(\alph*)]
\item $(1,1)$
\item $(2,2)$
\item $(1,2)$
\item $(2,1)$
\end{enumerate}

\item La pendiente de la recta que pasa por $(2,5)$ y $(4,9)$ es:
\begin{enumerate}[label=(\alph*)]
\item $2$
\item $1$
\item $3$
\item $4$
\end{enumerate}

\item La ecuación de la recta con pendiente $3$ que pasa por $(0,2)$ es:
\begin{enumerate}[label=(\alph*)]
\item $y = 3x + 2$
\item $y = 2x + 3$
\item $y = 3x - 2$
\item $y = 2x - 3$
\end{enumerate}

\item ¿Cuál es el centro de la circunferencia $(x-3)^2 + (y+2)^2 = 25$?
\begin{enumerate}[label=(\alph*)]
\item $(3,-2)$
\item $(-3,2)$
\item $(3,2)$
\item $(-3,-2)$
\end{enumerate}

\item El radio de la circunferencia $x^2 + y^2 = 16$ es:
\begin{enumerate}[label=(\alph*)]
\item $4$
\item $8$
\item $2$
\item $16$
\end{enumerate}

\item La ecuación de la recta que pasa por $(1,1)$ y es paralela a $y = 2x + 3$ es:
\begin{enumerate}[label=(\alph*)]
\item $y = 2x - 1$
\item $y = 2x + 1$
\item $y = -2x + 1$
\item $y = -2x - 1$
\end{enumerate}

\item El vértice de la parábola $y = x^2 - 4x + 3$ es:
\begin{enumerate}[label=(\alph*)]
\item $(2,-1)$
\item $(-2,1)$
\item $(2,1)$
\item $(-2,-1)$
\end{enumerate}

\item La distancia del punto $P(3,4)$ a la recta $x + y - 2 = 0$ es:
\begin{enumerate}[label=(\alph*)]
\item $\frac{5}{\sqrt{2}}$
\item $\frac{5}{2}$
\item $5$
\item $\sqrt{5}$
\end{enumerate}

\item El área del triángulo con vértices $(0,0)$, $(4,0)$ y $(0,3)$ es:
\begin{enumerate}[label=(\alph*)]
\item $6$
\item $12$
\item $8$
\item $10$
\end{enumerate}

\item La ecuación de la circunferencia con centro $(2,-1)$ y radio $3$ es:
\begin{enumerate}[label=(\alph*)]
\item $(x-2)^2 + (y+1)^2 = 9$
\item $(x+2)^2 + (y-1)^2 = 9$
\item $(x-2)^2 + (y+1)^2 = 3$
\item $(x+2)^2 + (y-1)^2 = 3$
\end{enumerate}

\item La pendiente de la recta perpendicular a $y = \frac{1}{2}x + 3$ es:
\begin{enumerate}[label=(\alph*)]
\item $-2$
\item $2$
\item $-\frac{1}{2}$
\item $\frac{1}{2}$
\end{enumerate}

\item El foco de la parábola $y^2 = 8x$ está en:
\begin{enumerate}[label=(\alph*)]
\item $(2,0)$
\item $(0,2)$
\item $(4,0)$
\item $(0,4)$
\end{enumerate}

\item Los focos de la elipse $\frac{x^2}{25} + \frac{y^2}{9} = 1$ son:
\begin{enumerate}[label=(\alph*)]
\item $(\pm 4,0)$
\item $(0,\pm 4)$
\item $(\pm 3,0)$
\item $(0,\pm 3)$
\end{enumerate}

\item La ecuación de la recta tangente a la circunferencia $x^2+y^2=25$ en el punto $(3,4)$ es:
\begin{enumerate}[label=(\alph*)]
\item $3x+4y=25$
\item $3x-4y=25$
\item $4x+3y=25$
\item $4x-3y=25$
\end{enumerate}

\item El valor de $k$ para que la recta $y = kx + 2$ sea tangente a la parábola $y = x^2$ es:
\begin{enumerate}[label=(\alph*)]
\item $k = \pm 2\sqrt{2}$
\item $k = \pm 2$
\item $k = \pm \sqrt{2}$
\item $k = \pm 4$
\end{enumerate}

\item La excentricidad de la elipse $\frac{x^2}{16} + \frac{y^2}{7} = 1$ es:
\begin{enumerate}[label=(\alph*)]
\item $\frac{3}{4}$
\item $\frac{4}{3}$
\item $\frac{1}{2}$
\item $\frac{\sqrt{7}}{4}$
\end{enumerate}

\item Las asíntotas de la hipérbola $\frac{x^2}{9} - \frac{y^2}{16} = 1$ son:
\begin{enumerate}[label=(\alph*)]
\item $y = \pm \frac{4}{3}x$
\item $y = \pm \frac{3}{4}x$
\item $y = \pm \frac{4}{3}y$
\item $y = \pm \frac{3}{4}y$
\end{enumerate}

\item El ángulo entre las rectas $x + y = 0$ y $x - y = 0$ es:
\begin{enumerate}[label=(\alph*)]
\item $90^\circ$
\item $45^\circ$
\item $60^\circ$
\item $30^\circ$
\end{enumerate}

\item La ecuación de la recta que pasa por la intersección de $x+y=3$ y $x-y=1$ y es perpendicular a $2x+y=0$ es:
\begin{enumerate}[label=(\alph*)]
\item $x-2y+1=0$
\item $x+2y-1=0$
\item $2x-y+1=0$
\item $2x+y-1=0$
\end{enumerate}

\item El lugar geométrico de los puntos equidistantes de $(2,0)$ y $(-2,0)$ es:
\begin{enumerate}[label=(\alph*)]
\item $x=0$
\item $y=0$
\item $x^2+y^2=4$
\item $y=x$
\end{enumerate}

\item La ecuación de la parábola con vértice en el origen y foco $(0,3)$ es:
\begin{enumerate}[label=(\alph*)]
\item $x^2 = 12y$
\item $y^2 = 12x$
\item $x^2 = 6y$
\item $y^2 = 6x$
\end{enumerate}

\item El centro de la hipérbola $9x^2 - 16y^2 = 144$ es:
\begin{enumerate}[label=(\alph*)]
\item $(0,0)$
\item $(4,0)$
\item $(0,3)$
\item $(3,4)$
\end{enumerate}

\item La ecuación de la recta que pasa por $(1,2)$ y forma un ángulo de $45^\circ$ con el eje $x$ es:
\begin{enumerate}[label=(\alph*)]
\item $y = x + 1$
\item $y = x + 2$
\item $y = -x + 3$
\item $y = -x + 1$
\end{enumerate}

\item La suma de las coordenadas del punto de intersección de las rectas $2x+3y=12$ y $x-2y=1$ es:
\begin{enumerate}[label=(\alph*)]
\item $5$
\item $4$
\item $6$
\item $7$
\end{enumerate}

\end{enumerate}
\end{multicols}
\end{document}